\chapter{Preliminares matem\'aticos}

\section{Introducci\'on}

Las ecuaciones diferenciales parciales (EDPs) constituyen una herramienta central para modelar fenómenos físicos, biológicos y de ingeniería. Dentro de la clase de ecuaciones elípticas, la ecuación de Poisson ocupa un lugar prototípico por su aparición en electrostática, conducción de calor y mecánica de fluidos, y porque permite estudiar con claridad cuestiones de existencia, regularidad, estabilidad y aproximación numérica.

Los métodos de diferencias finitas, elementos finitos y volúmenes finitos poseen una teoría de convergencia ampliamente desarrollada y son particularmente eficaces en problemas de dimensión baja o moderada \cite{strang1973, leveque2007, babuska1994}. Su aplicación requiere, sin embargo, una discretización del dominio y, según la geometría y la dimensión, puede demandar una construcción de malla o una estrategia de integración no trivial.

Las redes neuronales físicamente informadas (PINNs) ofrecen una formulación alternativa: representan la solución mediante una red neuronal y entrenan sus parámetros penalizando el residuo diferencial y el defecto en las condiciones de contorno. La metodología fue popularizada por Raissi, Perdikaris y Karniadakis \cite{raissi2019}. Su interés no reside en una superioridad universal frente a los métodos clásicos, sino en la posibilidad de combinar ecuaciones, condiciones de contorno y datos observacionales dentro de un mismo problema de optimización.

Esta tesis se concentra en una pregunta más específica: bajo qué hipótesis una pérdida PINN pequeña, evaluada en un conjunto finito de puntos de colocación, permite controlar el residuo continuo y el error funcional de la aproximación. La ecuación de Poisson se emplea como banco de pruebas porque dispone de una teoría elíptica bien establecida y de soluciones de referencia con las que es posible contrastar los resultados computacionales.

\section{EDPs elípticas lineales}

Sea $U\subset\mathbb{R}^d$ un dominio acotado y sea $\Gamma\subseteq\partial U$ una porción de la frontera. En el marco probabilístico del Capítulo~\ref{chap:convergencia} se supondrá, siguiendo a Shin et al., que $\Gamma$ es un subconjunto cerrado de $\partial U$. Consideramos el problema lineal
\[
\begin{cases}
    \mathcal{L}[u] = f & \text{en } U, \\
    \mathcal{B}[u] = g & \text{en } \Gamma,
\end{cases}
\]
donde $f:U\to\mathbb{R}$, $g:\Gamma\to\mathbb{R}$ y, para condiciones de Dirichlet,
$\mathcal{B}[u]=u|_{\Gamma}$. Cuando el dato se prescribe en toda la frontera se escribirá explícitamente
$\Gamma=\partial U$.

Siguiendo la formulación elíptica utilizada por Shin, Darbon y
Karniadakis~\cite{shin2020convergence}, escribimos el operador de modo
consistente con su expresión en forma no divergente:
\begin{equation}\label{eq:operador_eliptico_corregido}
\begin{aligned}
\mathcal{L}[u]
={}&\sum_{i,j=1}^{d}D_i\!\left(a^{ij}(x)D_j u\right)
 +\sum_{i=1}^{d}D_i\!\left(b^i(x)u\right)\\
 &+\sum_{i=1}^{d}c^i(x)D_i u+d(x)u.
\end{aligned}
\end{equation}
Aquí $D_i=\partial/\partial x_i$ y $d(x)$ es el coeficiente de orden
cero, que no debe confundirse con la dimensión espacial $d$.
La diferenciación mediante la regla del producto da
\[
\begin{aligned}
\sum_{i,j}D_i(a^{ij}D_j u)
&=\sum_{i,j}a^{ij}D_{ij}u
 +\sum_j\!\left(\sum_i D_i a^{ij}\right)D_j u,\\
\sum_i D_i(b^i u)&=\sum_i b^iD_i u
 +\left(\sum_iD_i b^i\right)u.
\end{aligned}
\]
En consecuencia, si $a^{ij},b^i\in C^1(U)$, la forma no divergente es
\begin{equation}\label{eq:operador_eliptico_no_divergente}
\mathcal{L}[u]
=\sum_{i,j=1}^{d}a^{ij}(x)D_{ij}u
 +\sum_{j=1}^{d}\widetilde b^{\,j}(x)D_j u
 +\widetilde c(x)u,
\end{equation}
con
\[
\widetilde b^{\,j}=\sum_{i=1}^{d}D_i a^{ij}+b^j+c^j,
\qquad
\widetilde c=\sum_{i=1}^{d}D_i b^i+d.
\]
La expresión \eqref{eq:operador_eliptico_corregido} fija la convención
de índices empleada en esta tesis y es la que se utilizará en la
Hipótesis~3 y en las estimaciones de estabilidad. Se trata de una
precisión de notación, no de una contribución matemática reivindicada.
En particular, el principio del máximo exige la condición de signo
$\widetilde c\leq0$ cuando se aplique a este operador. Las estimaciones de
Schauder requieren además regularidad suficiente de $\widetilde b$ y
$\widetilde c$ \cite{gilbarg2001elliptic}.

En los experimentos de Poisson se toma $\mathcal L=\Delta$; cuando se usa
$-\Delta$ por conveniencia en alguna formulación numérica, el término
fuente se cambia de signo de manera explícita. Esta convención no modifica
las demostraciones abstractas del presente capítulo.

El operador es uniformemente elíptico si existe $\lambda_0>0$ tal que
\[
\sum_{i,j=1}^{d}a^{ij}(x)\xi_i\xi_j
\geq
\lambda_0|\xi|^2,
\qquad
\forall\xi\in\mathbb{R}^d,\quad \forall x\in U.
\]

\section{Espacios de funciones}

Las estimaciones que emplearemos para analizar la convergencia de las PINNs
involucran tres tipos de normas: las de H\"older, las de integrabilidad
cuadr\'atica y las de Sobolev. Introducimos primero la regularidad de H\"older,
porque permite relacionar la informaci\'on disponible en los puntos de
entrenamiento con el comportamiento de una funci\'on en el dominio completo.

\subsection*{Espacios de H\"older}

Sea $U\subset\mathbb R^d$ un dominio acotado y $0<\alpha\leq 1$. Para una
funci\'on $v:U\to\mathbb R$, consideramos la seminorma
\[
[v]_{\alpha;U}
:=\sup_{\substack{x,y\in U\\x\neq y}}
\frac{|v(x)-v(y)|}{\|x-y\|^\alpha}.
\]
Su finitud expresa que las variaciones de $v$ est\'an controladas por una
potencia de la distancia:
\[
|v(x)-v(y)|\leq[v]_{\alpha;U}\,\|x-y\|^\alpha,
\qquad x,y\in U.
\]
El caso $\alpha=1$ corresponde a la continuidad Lipschitz; para
$0<\alpha<1$ se obtiene una condici\'on de regularidad m\'as d\'ebil.

Para extender esta noci\'on a las derivadas, utilizamos multi\'indices
$\beta=(\beta_1,\ldots,\beta_d)\in\mathbb N_0^d$, con
$|\beta|=\beta_1+\cdots+\beta_d$, y escribimos
\[
D^\beta u
=\frac{\partial^{|\beta|}u}
{\partial x_1^{\beta_1}\cdots\partial x_d^{\beta_d}}.
\]
Si $u\in C^k(U)$, definimos, para $0\leq j\leq k$,
\[
[u]_{j,0;U}:=
\max_{|\beta|=j}\sup_{x\in U}|D^\beta u(x)|,
\qquad
[u]_{k,\alpha;U}:=
\max_{|\beta|=k}[D^\beta u]_{\alpha;U}.
\]
El espacio $C^{k,\alpha}(\overline U)$ est\'a formado por las funciones de
$C^k(\overline U)$ cuyas derivadas de orden $k$ son uniformemente
H\"older continuas. Lo equipamos con la norma
\[
\|u\|_{C^{k,\alpha}(\overline U)}
:=\sum_{j=0}^{k}[u]_{j,0;U}+[u]_{k,\alpha;U}.
\]
En particular, para $k=0$ escribimos tambi\'en $C^\alpha(\overline U)$.
Con la norma anterior, $C^{k,\alpha}(\overline U)$ es un espacio de
Banach. La notaci\'on $C^{k,\alpha}_{\mathrm{loc}}(U)$ se reservar\'a
para funciones que poseen esta regularidad sobre cada compacto contenido
en $U$, sin exigir el mismo control hasta la frontera.

Esta distinci\'on tendr\'a un papel concreto en el Cap\'itulo III. Para una
aproximaci\'on $h$ de la soluci\'on de $\mathcal L[u]=f$,
sea $r_h=\mathcal L[h]-f$ su residuo. Si $x_i$ es un punto de colocaci\'on,
entonces
\[
|r_h(x)|\leq |r_h(x_i)|+
\bigl([\mathcal L[h]]_{\alpha;U}+[f]_{\alpha;U}\bigr)
\|x-x_i\|^\alpha,
\qquad x\in U.
\]
La desigualdad permite distinguir el residuo evaluado en los puntos de
entrenamiento del control que se necesita entre esos puntos. En este
sentido se emplean las seminormas de H\"older en la regularizaci\'on y
en las cotas de generalizaci\'on estudiadas posteriormente.
La notaci\'on y su aplicaci\'on a PINNs se relacionan con la secci\'on 2.1
de Shin et al.~\cite{shin2020convergence}; el marco cl\'asico de regularidad
el\'iptica puede consultarse en \cite{gilbarg2001elliptic}.

\subsection*{Espacios \(L^2\)}

Sea \(U \subseteq \mathbb{R}^d\) un dominio acotado. El espacio \(L^2(U)\) se define como el conjunto de funciones reales medibles \(f: U \rightarrow \mathbb{R}\) tales que el cuadrado de su valor absoluto es integrable en \(U\). Más formalmente,

\[
L^2(U) = \left\{ f : U \rightarrow \mathbb{R} \,\Big|\, \int_U |f(x)|^2\,dx < \infty \right\}.
\]

El espacio \(L^2(U)\) es un espacio de Hilbert equipado con el producto interno dado por:

\[
\langle f,g \rangle_{L^2(U)} = \int_U f(x)g(x)\,dx, \quad f,g \in L^2(U),
\]

y con la norma inducida correspondiente definida por:

\[
\|f\|_{L^2(U)} = \left(\int_U |f(x)|^2\,dx\right)^{1/2}.
\]

Esta estructura dota al espacio \(L^2(U)\) de propiedades útiles, incluyendo la completitud respecto a esta norma, lo que resulta esencial para diversos análisis matemáticos, particularmente en el estudio de ecuaciones diferenciales parciales y teoría de aproximación.

\subsection*{Espacios de Sobolev}

Sea $U\subset\mathbb{R}^d$ acotado y de frontera Lipschitz. Para $k\in\mathbb{N}$ definimos

\[
H^k(U)=\big\{u\in L^2(U): D^\beta u\in L^2(U)\ \text{para todo multiíndice }|\beta|\le k\big\},
\]

con norma

\[
\|u\|_{H^k(U)}^2=\sum_{|\beta|\le k}\|D^\beta u\|_{L^2(U)}^2.
\]

En particular,

\[
H^1(U)=\{u\in L^2(U): \nabla u\in (L^2(U))^d\},\qquad
\|u\|_{H^1(U)}^2=\|u\|_{L^2(U)}^2+\|\nabla u\|_{L^2(U)}^2.
\]

Definimos $H_0^1(U)$ como el cierre de $C_c^\infty(U)$ en $H^1(U)$; para dominios Lipschitz, coincide con el subespacio de funciones cuya traza se anula en toda $\partial U$. Cuando solo se prescribe traza nula en una porción $\Gamma\subseteq\partial U$, emplearemos la notación $H^1_{0,\Gamma}(U)$.

\section{Principio del máximo}

El principio del máximo proporciona unicidad y estimaciones \emph{a priori} para problemas elípticos. En una formulación típica, si $\mathcal{L}$ es uniformemente elíptico, el coeficiente de orden cero satisface la condición de signo requerida y $\mathcal{L}[u]\geq 0$ en $U$, entonces
\[
\max_{\overline U}u\leq \max\bigl\{0,\max_{\partial U}u\bigr\}.
\]
En particular, cuando $u$ es armónica o subarmónica bajo las hipótesis usuales, el máximo se controla por el dato prescrito en la frontera completa $\partial U$. No debe sustituirse $\partial U$ por una porción propia $\Gamma$ sin imponer condiciones adicionales sobre el resto de la frontera. Véase \cite{gilbarg2001elliptic}.

En esta tesis, el principio del máximo se utiliza no solo como herramienta teórica clásica, sino como mecanismo para \emph{acotar el error de una solución aproximada}. Sea $u$ la solución exacta del problema de Poisson y sea $h_\theta$ la aproximación obtenida mediante una PINN. Definimos el error
\[
e := u - h_\theta.
\]
Entonces $e$ satisface un problema elíptico con datos ``pequeños'' dados por el error en la frontera y el residuo de la EDP. De manera esquemática, podemos escribir
\[
L[e] = f - L[h_\theta] \quad \text{en } U, \qquad
e = g - h_\theta \quad \text{en } \partial U,
\]
donde $f$ y $g$ son los datos exactos del problema.

El principio del máximo permite mostrar que, bajo hipótesis de elipticidad uniforme y regularidad adecuadas, el error global $\|e\|_{C^0(\overline{U})}$ está controlado por:
\begin{itemize}
  \item qué tan bien la PINN satisface las condiciones de frontera, es decir, el tamaño de $g-h_\theta$ en $\partial U$, y
  \item qué tan bien minimiza el residuo de la EDP en el dominio, es decir, el tamaño de $f - L[h_\theta]$ en $U$.
\end{itemize}
En particular, si
\begin{enumerate}
  \item la PINN aproxima bien las condiciones de frontera (el error $|g-h_\theta|$ en $\partial U$ es pequeño), y
  \item el residuo $|L[h_\theta] - f|$ es pequeño en $U$ (la física se respeta en el interior),
\end{enumerate}
entonces el principio del máximo, junto con estimaciones elípticas estándar, permite concluir que el \emph{error global} $|u - h_\theta|$ está controlado en todo el dominio. Esta idea es central en los resultados de convergencia que se desarrollan en el Capítulo~III, donde se formaliza el paso desde la minimización de la pérdida PINN (residuo + error de frontera) hasta estimaciones de convergencia en normas $C^0$ y $H^1$.


\section{Celdas de Voronoi}

Sea \( U \subseteq \mathbb{R}^d \) un dominio acotado y considérese un conjunto finito de puntos distintos \( X = \{ x_i \}_{i=1}^{n} \subset U \). A cada punto \( x_i \) del conjunto \( X \) se le asocia una región denominada \emph{celda de Voronoi}, que se define formalmente a continuación.

\textbf{Definición.} La celda de Voronoi asociada a un punto \( x_i \in X \), denotada por \( V_i \), se define como el conjunto de puntos del dominio \( U \) que se encuentran más cerca del punto \( x_i \) que de cualquier otro punto del conjunto \( X \). Matemáticamente, esta definición se expresa como:
\[
V_i := \{ x \in U : \| x - x_i \| \leq \| x - x_j \|, \; \forall j \neq i, j=1,\dots,n \}.
\]

Las celdas cubren $U$, pero con la definición mediante ``$\leq$'' pueden compartir puntos equidistantes. Para obtener una partición medible se fija una regla determinista de desempate. Con esa convención,
\[
U=\bigcup_{i=1}^{n}V_i,
\qquad
V_i\cap V_j=\varnothing\quad(i\neq j),
\]
salvo, equivalentemente, conjuntos de frontera de medida nula.

\begin{figure}[h!]
    \centering
    \includegraphics[width=0.43\linewidth]{imagenes/voronoi.png}
    \caption{Celdas de Voronoi en un dominio bidimensional.}
    \label{fig:voronoi}
\end{figure}

Si \(\mu\) es una medida de probabilidad en \(U\), la medida de una celda de Voronoi se define como

\[
\mu(V_i) = \int_{V_i} d\mu(x).
\]

Para una muestra empírica \(T_{m} = \{ z_k \}_{k=1}^{m} \subset U\), la medida empírica de \(V_i\) es

\[
\mu_{m}(V_i) = \frac{1}{m} \sum_{k=1}^{m} \delta_{z_k}(V_i),
\]

donde \(\delta_{z_k}(V_i)=1\) si \(z_k \in V_i\) y \(0\) en caso contrario.

\section{Existencia y regularidad de Schauder}

\textbf{Teorema.}
Sea $U\subset\mathbb{R}^d$ un dominio acotado con frontera $C^{2,\alpha}$, $0<\alpha<1$, y sea
\[
\mathcal{L}u
=
\sum_{i,j=1}^{d}a^{ij}(x)D_{ij}u
+
\sum_{i=1}^{d}b^i(x)D_i u
+
c(x)u
\]
un operador uniformemente elíptico con $a^{ij}=a^{ji}$ y coeficientes
$a^{ij},b^i,c\in C^{0,\alpha}(\overline U)$. Supóngase además que la condición de signo necesaria para unicidad, por ejemplo $c\leq 0$, se cumple. Entonces, para todo
$f\in C^{0,\alpha}(\overline U)$ y $g\in C^{2,\alpha}(\partial U)$, el problema
\[
\begin{cases}
\mathcal{L}u=f & \text{en }U,\\
u=g & \text{en }\partial U
\end{cases}
\]
admite una única solución $u\in C^{2,\alpha}(\overline U)$. Además, existe una constante $C$, dependiente del dominio, de la elipticidad y de las normas de los coeficientes, tal que
\[
\|u\|_{C^{2,\alpha}(\overline U)}
\leq
C\Bigl(
\|f\|_{C^{0,\alpha}(\overline U)}
+
\|g\|_{C^{2,\alpha}(\partial U)}
\Bigr).
\]

\paragraph{Comentario.}
La demostración estándar combina estimaciones de Schauder, el principio del máximo y el método de continuidad; remitimos a \cite{gilbarg2001elliptic} para el enunciado preciso y su prueba. Este es un resultado de \emph{regularidad global} para fronteras suaves.
No se aplica literalmente al cuadrado $(-1,1)^2$, cuya frontera tiene esquinas.
En el Capítulo~III, para el problema de Dirichlet en dominios que admiten
esferas exteriores, se emplea en cambio un resultado de existencia y
regularidad \emph{interior} con continuidad hasta la frontera; no se
identifican las hipótesis de ambos teoremas.


